% GENERATED FILE. Edit canonical lessons, then run npm run generate:books.
% canonical-source-hash: fnv1a64:0bdd784a9801e3b0
% canonical-lessons: SA-C434-uttamarnah, SA-C434-adhamarnah, SA-C434-kusidam, SA-C434-niksepah, SA-C434-niryatayati

\chapter{Loans and Deposits}
\label{ch:sa-loans-and-deposits}
\hlchaptermodality{434}

\begin{chapteropening}
\textbf{By the end of this chapter:} \emph{I can talk about lending, borrowing and paying back.}

Complete the last lesson of chapter 434: I can talk about lending, borrowing and paying back.
\end{chapteropening}

\section[\texorpdfstring{\sk{उत्तमर्णः} (uttamarṇaḥ)}{उत्तमर्णः (uttamarṇaḥ)}]{\sk{उत्तमर्णः} (uttamarṇaḥ) — a creditor}

\label{lesson:SA-C434-uttamarnah}



\begin{quote}
Before the new one: say the Sanskrit for a share, then the Sanskrit for a guarantor.
\end{quote}

\subsection*{You'll want to know: \sk{उत्तमर्णः}}
\textbf{\sk{उत्तमर्णः}} — \emph{uttamarṇaḥ} — \textquotedblleft{}a creditor\textquotedblright{}.

\begin{quote}
\sk{उत्तमर्णः} \sk{ऋणं} \sk{ददाति।}
\end{quote}

\begin{quote}
\emph{uttamarṇaḥ ṛṇaṃ dadāti.}
\end{quote}

\begin{quote}
The creditor gives a loan.
\end{quote}

A masculine noun.

\begin{cousinweb}
\textbf{\sk{उत्तमर्णः}} (\emph{uttamarṇaḥ}) is built from \emph{uttama}, \textquotedblleft{}highest, upper\textquotedblright{}, and \emph{ṛṇa}, \textquotedblleft{}a debt\textquotedblright{}; MW gives the compound as a creditor (MW).
\end{cousinweb}

\begin{grammarlens}[title={What you've built}]
One more word for getting around.
\end{grammarlens}

\subsection*{Your turn}
\begin{itemize}
\raggedright
  \item \emph{Say it:} \emph{uttamarṇaḥ}
  \item \emph{Say it:} \emph{uttamarṇaḥ}, once more
  \item \emph{Say it:} \emph{uttamarṇaḥ ṛṇaṃ dadāti.} — The creditor gives a loan.
  \item \emph{Recall:} say the Sanskrit for a share, then the Sanskrit for a guarantor, then say \emph{uttamarṇaḥ} again
\end{itemize}

\begin{tcolorbox}[breakable,colback=teal!4,colframe=teal!35!black,title={Before you move on}]
What does \sk{उत्तमर्णः} mean? (\textquotedblleft{}A creditor\textquotedblright{}.) Say it once more.
\end{tcolorbox}
\section[\texorpdfstring{\sk{अधमर्णः} (adhamarṇaḥ)}{अधमर्णः (adhamarṇaḥ)}]{\sk{अधमर्णः} (adhamarṇaḥ) — a debtor}

\label{lesson:SA-C434-adhamarnah}



\begin{quote}
Before the new one: say the Sanskrit for a guarantor, then the Sanskrit for a creditor.
\end{quote}

\subsection*{You'll want to know: \sk{अधमर्णः}}
\textbf{\sk{अधमर्णः}} — \emph{adhamarṇaḥ} — \textquotedblleft{}a debtor\textquotedblright{}.

\begin{quote}
\sk{अधमर्णः} \sk{शनैः} \sk{ऋणं} \sk{शोधयति।}
\end{quote}

\begin{quote}
\emph{adhamarṇaḥ śanaiḥ ṛṇaṃ śodhayati.}
\end{quote}

\begin{quote}
The debtor slowly clears the debt.
\end{quote}

A masculine noun.

\begin{cousinweb}
\textbf{\sk{अधमर्णः}} (\emph{adhamarṇaḥ}) is built from \emph{adhama}, \textquotedblleft{}lowest\textquotedblright{}, and \emph{ṛṇa}, \textquotedblleft{}a debt\textquotedblright{}; MW gives the compound as a debtor (MW).
\end{cousinweb}

\begin{grammarlens}[title={What you've built}]
One more word for getting around.
\end{grammarlens}

\subsection*{Your turn}
\begin{itemize}
\raggedright
  \item \emph{Say it:} \emph{adhamarṇaḥ}
  \item \emph{Say it:} \emph{adhamarṇaḥ}, once more
  \item \emph{Say it:} \emph{adhamarṇaḥ śanaiḥ ṛṇaṃ śodhayati.} — The debtor slowly clears the debt.
  \item \emph{Recall:} say the Sanskrit for a guarantor, then the Sanskrit for a creditor, then say \emph{adhamarṇaḥ} again
\end{itemize}

\begin{tcolorbox}[breakable,colback=teal!4,colframe=teal!35!black,title={Before you move on}]
What does \sk{अधमर्णः} mean? (\textquotedblleft{}A debtor\textquotedblright{}.) Say it once more.
\end{tcolorbox}
\section[\texorpdfstring{\sk{कुसीदम्} (kusīdam)}{कुसीदम् (kusīdam)}]{\sk{कुसीदम्} (kusīdam) — money lent at interest}

\label{lesson:SA-C434-kusidam}



\begin{quote}
Before the new one: say the Sanskrit for a creditor, then the Sanskrit for a debtor.
\end{quote}

\subsection*{You'll want to know: \sk{कुसीदम्}}
\textbf{\sk{कुसीदम्}} — \emph{kusīdam} — \textquotedblleft{}money lent at interest\textquotedblright{}.

\begin{quote}
\sk{कुसीदेन} \sk{धनिकस्य} \sk{धनं} \sk{वर्धते।}
\end{quote}

\begin{quote}
\emph{kusīdena dhanikasya dhanaṃ vardhate.}
\end{quote}

\begin{quote}
The rich man's wealth grows by lending at interest.
\end{quote}

A neuter noun.

\begin{cousinweb}
\textbf{\sk{कुसीदम्}} (\emph{kusīdam}) is the word MW gives for a loan to be repaid with interest and for lending money upon interest; MW suggests \emph{ku} and the root √sad with a query, so its origin is uncertain (MW).
\end{cousinweb}

\begin{grammarlens}[title={What you've built}]
One more word for getting around.
\end{grammarlens}

\subsection*{Your turn}
\begin{itemize}
\raggedright
  \item \emph{Say it:} \emph{kusīdam}
  \item \emph{Say it:} \emph{kusīdam}, once more
  \item \emph{Say it:} \emph{kusīdena dhanikasya dhanaṃ vardhate.} — The rich man's wealth grows by lending at interest.
  \item \emph{Recall:} say the Sanskrit for a creditor, then the Sanskrit for a debtor, then say \emph{kusīdam} again
\end{itemize}

\begin{tcolorbox}[breakable,colback=teal!4,colframe=teal!35!black,title={Before you move on}]
What does \sk{कुसीदम्} mean? (\textquotedblleft{}Money lent at interest\textquotedblright{}.) Say it once more.
\end{tcolorbox}
\section[\texorpdfstring{\sk{निक्षेपः} (nikṣepaḥ)}{निक्षेपः (nikṣepaḥ)}]{\sk{निक्षेपः} (nikṣepaḥ) — a deposit}

\label{lesson:SA-C434-niksepah}



\begin{quote}
Before the new one: say the Sanskrit for a debtor, then the Sanskrit for money lent at interest.
\end{quote}

\subsection*{You'll want to know: \sk{निक्षेपः}}
\textbf{\sk{निक्षेपः}} — \emph{nikṣepaḥ} — \textquotedblleft{}a deposit\textquotedblright{}.

\begin{quote}
\sk{सः} \sk{वित्तकोषे} \sk{निक्षेपं} \sk{स्थापयति।}
\end{quote}

\begin{quote}
\emph{saḥ vittakoṣe nikṣepaṃ sthāpayati.}
\end{quote}

\begin{quote}
He places a deposit in the bank.
\end{quote}

A masculine noun.

\begin{cousinweb}
\textbf{\sk{निक्षेपः}} (\emph{nikṣepaḥ}) is built from \emph{ni}, \textquotedblleft{}down\textquotedblright{}, and the root √kṣip, \textquotedblleft{}to throw\textquotedblright{}, which also gives \emph{nikṣipati}, \textquotedblleft{}puts in, deposits\textquotedblright{}; MW gives \emph{nikṣepa} as putting down and as a deposit, a pledge (MW).
\end{cousinweb}

\begin{grammarlens}[title={What you've built}]
One more word for getting around.
\end{grammarlens}

\subsection*{Your turn}
\begin{itemize}
\raggedright
  \item \emph{Say it:} \emph{nikṣepaḥ}
  \item \emph{Say it:} \emph{nikṣepaḥ}, once more
  \item \emph{Say it:} \emph{saḥ vittakoṣe nikṣepaṃ sthāpayati.} — He places a deposit in the bank.
  \item \emph{Recall:} say the Sanskrit for a debtor, then the Sanskrit for money lent at interest, then say \emph{nikṣepaḥ} again
\end{itemize}

\begin{tcolorbox}[breakable,colback=teal!4,colframe=teal!35!black,title={Before you move on}]
What does \sk{निक्षेपः} mean? (\textquotedblleft{}A deposit\textquotedblright{}.) Say it once more.
\end{tcolorbox}
\section[\texorpdfstring{\sk{निर्यातयति} (niryātayati)}{निर्यातयति (niryātayati)}]{\sk{निर्यातयति} (niryātayati) — pays back; returns}

\label{lesson:SA-C434-niryatayati}



\begin{quote}
Before the new one: say the Sanskrit for money lent at interest, then the Sanskrit for a deposit.
\end{quote}

\subsection*{You'll want to know: \sk{निर्यातयति}}
\textbf{\sk{निर्यातयति}} — \emph{niryātayati} — \textquotedblleft{}pays back; returns\textquotedblright{}.

\begin{quote}
\sk{अधमर्णः} \sk{उत्तमर्णाय} \sk{धनं} \sk{निर्यातयति।}
\end{quote}

\begin{quote}
\emph{adhamarṇaḥ uttamarṇāya dhanaṃ niryātayati.}
\end{quote}

\begin{quote}
The debtor pays the money back to the creditor.
\end{quote}

Past (he …): \sk{सः} \sk{धनं} \sk{निर्यातितवान्}, \emph{saḥ dhanaṃ niryātitavān}.

\begin{cousinweb}
\textbf{\sk{निर्यातयति}} (\emph{niryātayati}) is built from \emph{nir}, \textquotedblleft{}out\textquotedblright{}, and the causative of the root √yat, which also gives \emph{yatate}, \textquotedblleft{}tries\textquotedblright{}; MW gives its participle \emph{niryātita} as restored, returned, requited (MW).
\end{cousinweb}

\begin{grammarlens}[title={What you've built}]
One more word for getting around.
\end{grammarlens}

\subsection*{Your turn}
\begin{itemize}
\raggedright
  \item \emph{Say it:} \emph{niryātayati}
  \item \emph{Say it:} \emph{niryātayati}, once more
  \item \emph{Say it:} \emph{adhamarṇaḥ uttamarṇāya dhanaṃ niryātayati.} — The debtor pays the money back to the creditor.
  \item \emph{Recall:} say the Sanskrit for money lent at interest, then the Sanskrit for a deposit, then say \emph{niryātayati} again
\end{itemize}

\begin{tcolorbox}[breakable,colback=teal!4,colframe=teal!35!black,title={Before you move on}]
What does \sk{निर्यातयति} mean? (\textquotedblleft{}Pays back; returns\textquotedblright{}.) Say it once more.
\end{tcolorbox}
